\section{Experiments}
\label{sec:experiments}
\subsection{Datasets, metrics and comparison protocol}
We use one fixed provider recipe across 31 scene jobs. Four LERF scenes yield
two outputs, so the cohort comprises five evaluation tasks rather than five
independent datasets. All reported numbers are percentages. The existing
benchmark rooms have been used during development; they are not an untouched
final test. Scene-level intervals describe this cohort rather than certify
generalization after repeated research feedback.

\begin{table}[ht]
\centering\small
\begin{tabular}{lp{.25\linewidth}p{.43\linewidth}}
\toprule
Task & Cohort and primary metric & Input and evaluation contract \\
\midrule
LERF-2D & 4 scenes, 208 observations; mIoU and localization & Source-assisted text; semantic response rendered before thresholding. \\
LERF-3D selection & Same scenes; projection mIoU and Acc.@0.25 & Independent selected Gaussians rendered to 2D GT; not volumetric GT IoU. \\
ScanNet OVS & Paper8 Gaussian mIoU/mAcc; internal10 mesh mIoU & Fixed 19/15/10 menus; Gaussian and mesh domains reported separately. \\
NVOS & 8 tasks, 7 scene clusters; IoU & Retained positive-reference, cached-source variant; external prompt/RGB access declared. \\
UQIS & 9 rooms, 67 instances, 268 queries; set IoU and AP & Text, image, 2D point and 3D point; map-only query-time access. \\
\bottomrule
\end{tabular}
\caption{Evaluation units. LERF and UQIS average within scenes then across
scenes; NVOS averages eight tasks, with three seeds per task for LUDVIG.}
\label{tab:protocols}
\end{table}

\paragraph{Published results and verified reproductions.}
We reuse values reported in VALA Tables 1--2 and LUDVIG Tables 1/D, including
their cited baseline results~\cite{vala,ludvig}. No new external baseline training
is needed. Every local reproduction is linked to its frozen result hashes;
four-scene, eight-scene and eight-task aggregates are recomputed from the
original reports. Table origins are \textbf{P}: published, \textbf{R}: locally
reproduced, \textbf{B}: frozen provider and \textbf{M}: measured shared replacement.
An empty mechanism result is marked pending, never filled with a provider score.
We do not bold a global winner across different input contracts.

\subsection{Implementation and independent supervision}
We keep the 22 fit, three calibration and four development rooms, disjoint from
benchmark room identities. The inherited membership provider also used the fit
rooms. Natural-language episodes provide 629/207/175 fit/calibration/development
expressions. Original RGB-D and official instance annotations construct 290
natural image queries, split 237/28/25. Annotation geometry selects reference
views; neither teacher predictions nor predicted utility choose the query crop.
Reference frames, temporal neighbors and nearby poses are jointly excluded from
image candidate construction. At least half of mapping views and two
target-observing views remain. The full-scan geometry remains available; target
masks are not inputs to the query encoders. Original SigLIP and DINO image
evidence and complete decoders are preserved.

Category supervision adds 116 whole-menu episodes: one full map and three
chronological view folds per room. Only 88 fit-room episodes enter training;
calibration and development use full maps. SigLIP2, Mask-Adapter and literal
category SAM3 evidence are retained. View folds are correlated and do not add
independent rooms. Six classes are absent from calibration. One margin protects
text and image IoU and the three category outputs in every calibration room
(15 cells). Calibration and development feedback has informed model design.
No modality or benchmark identifier is an estimator input.

The latest evaluated parameterization splits the common potential into a
structural prior and a query-conditioned residual. Five room-wise folds supply
out-of-room prior predictions for residual targets; each branch uses 64 trees.
Permuted queries receive zero residual targets, an assumption tested by the
matched no-null control. The deployed comparison uses the same candidate pool,
one checkpoint and one margin (0.15052), selected before benchmark prediction.
All 19 ScanNet/UQIS scene-job predictions are sealed before scoring; only compact
signatures, models and reports persist. Dense predictions remain transient.

\subsection{Comparison with existing methods}
\paragraph{LERF semantic segmentation and projected selection.}
Tables~\ref{tab:lerf2d}--\ref{tab:lerf3d} reuse the published context and verified
Occam/VALA reproductions. The VALA reproduction follows its semantic and
evaluation path on compatible Occam RGB geometry, rather than retraining the
entire original geometry pipeline. The frozen v10 values quantify preserved
provider capability; the new shared decision has not yet been measured on these
two outputs. In particular, semantic-response improvements cannot be substituted
for independent selected-set improvements.
Our provider includes SAM3 text-conditioned source evidence and native SigLIP2
matching; external VALA scores are not matched-teacher controls. An independent
CPU audit exactly reproduces all 208 retained selected-set mask IoUs. Semantic
IoU/localization records are reaggregated, without a fresh semantic render in
that audit.

\begin{table}[ht]
\centering\small\begin{tabular}{llrr}
\toprule
Method & Origin & mIoU & Loc. acc. \\
\midrule
LERF \cite{vala} & P & 37.40 & 73.60 \\
LangSplat \cite{vala} & P & 51.40 & 84.30 \\
GAGS \cite{vala} & P & 54.10 & 81.70 \\
LangSplatV2 \cite{vala} & P & 59.90 & 84.10 \\
Occam's LGS \cite{vala} & P & 61.30 & 82.50 \\
VALA \cite{vala} & P & 61.70 & 86.40 \\
Occam's LGS & R & 63.62 & 82.85 \\
Ours (v10) & O & 70.40 & 92.42 \\
\bottomrule
\end{tabular}

\caption{LERF-2D. P values are reported in VALA Table 1~\cite{vala}.
Localization is peak-in-box accuracy, not ranking AP. B uses source-assisted
full teachers; it is not a gain of the proposed replacement decision.}
\label{tab:lerf2d}
\end{table}
\begin{table}[ht]
\centering\small\begin{tabular}{llrr}
\toprule
Method & Origin & mIoU & Acc.@0.25 \\
\midrule
OpenGaussian \cite{vala} & P & 38.36 & 51.43 \\
Dr. Splat \cite{vala} & P & 43.29 & 64.30 \\
InstanceGaussian \cite{vala} & P & 43.87 & 61.09 \\
CAGS \cite{vala} & P & 50.79 & 69.62 \\
VoteSplat \cite{vala} & P & 50.10 & 67.38 \\
Occam's LGS \cite{vala} & P & 47.22 & 74.84 \\
VALA \cite{vala} & P & 58.02 & 82.85 \\
VALA (compatible geometry) & R & 54.12 & 79.35 \\
Ours (v10) & O & 66.96 & 90.02 \\
\bottomrule
\end{tabular}

\caption{LERF selected-Gaussian projection. P values are reported in VALA
Table 1~\cite{vala}; Acc.@0.25 counts projected masks with IoU above 0.25.
VALA's published value is 58.02/82.85; 43.29/64.30 belongs to Dr. Splat.
The local VALA R row uses compatible RGB geometry.}
\label{tab:lerf3d}
\end{table}

\paragraph{ScanNet open-vocabulary segmentation.}
Table~\ref{tab:scannet} separates the paper-eight Gaussian metric from the
original ten-room mesh metric. For the metric bridge, we score unchanged v10
labels with the existing VALA evaluator: anisotropic Mahalanobis pseudo labels,
five-scale radius, top-1000 class-balanced votes and opacity--volume weights.
All eight scene predictions are hashed before opening their target labels.
This aligns the metric and scene list, but not the input carrier: v10 uses
provided unlabelled mesh geometry, while VALA uses optimized RGB Gaussians.
The bridge is a protocol diagnostic, not an additional learned gain. The
replacement policy retains every original menu on the internal ten-room track.
Independent scoring reproduces all ten mesh scenes and all eight Gaussian
scenes. All 1,134,207 pseudo labels agree, including 256 checks with the released
VALA code. On the same eight rooms, 19-class mesh mIoU is 61.0567 versus
66.9480 for Gaussian scoring: this 5.8913-point change is metric-dependent.

\begin{table}[ht]
\centering\small
\resizebox{\linewidth}{!}{\begin{tabular}{llrrrrrr}
\toprule
Method & Origin & IoU19 & Acc19 & IoU15 & Acc15 & IoU10 & Acc10 \\
\midrule
\midrule\multicolumn{8}{l}{\textit{Eight scenes: Gaussian scoring (input differences disclosed)}} \\
LangSplat \cite{vala} & P & 2.45 & 8.59 & 3.45 & 13.21 & 6.48 & 21.89 \\
OpenGaussian \cite{vala} & P & 27.73 & 42.01 & 29.67 & 46.15 & 39.93 & 57.34 \\
Dr. Splat \cite{vala} & P & 29.31 & 47.68 & 33.25 & 54.33 & 44.19 & 65.19 \\
Occam's LGS \cite{vala} & P & 31.93 & 48.93 & 34.25 & 53.71 & 45.16 & 64.39 \\
VALA \cite{vala} & P & 32.11 & 50.05 & 35.10 & 54.77 & 46.21 & 65.61 \\
VALA & R & 34.53 & 51.59 & 37.96 & 56.77 & 47.36 & 67.47 \\
Ours (v10, metric bridge) & O & 66.95 & 77.94 & 67.86 & 79.61 & 78.07 & 87.89 \\
\midrule\multicolumn{8}{l}{\textit{Ten scenes: mesh scoring (separate evaluation track)}} \\
Ours (v10) (10-mesh) & O & 61.01 & -- & 62.53 & -- & 72.49 & -- \\
Ours + replacement (L11) (10-mesh) & M & 61.01 & -- & 62.53 & -- & 72.49 & -- \\
\bottomrule
\end{tabular}
}
\caption{ScanNet OVS. P rows from VALA Table 2~\cite{vala}, R and the metric
bridge use paper8 Gaussian scoring. Rows marked 10-mesh use the original
ten-room mesh evaluation and cannot be ranked against the paper8 rows.
The metric bridge still differs in geometry and teacher access.}
\label{tab:scannet}
\end{table}

\paragraph{NVOS.}
Table~\ref{tab:nvos} includes the complete three-seed LUDVIG-SAM reproduction.
Its positive scribbles, all-view RGB geometry and online target-RGB access
follow the released implementation. Our retained provider uses a
positive-derived reference SAM mask and cached source tracks, without online
target RGB. These are contextual method comparisons with declared access
differences. A matched-input test is needed before attributing their difference
to the replacement mechanism.
The audit also finds a GT boundary convention difference: our retained scorer
uses $g>127$, while LUDVIG uses integer division by the mask maximum. Applying
only its GT convention changes our score from 86.5134 to 86.6835; historical
values remain unchanged. Switching bilinear implementations changes less than
0.00002 points. Borrowing its min--max and fixed 0.25 response threshold yields
86.6085 on our GT convention, and is not a full LUDVIG reproduction. Flower,
trex and orchids account for 78.04\% of the gap to the local 91.2577 reproduction.
Flower/orchids primarily miss foreground; trex also activates background.
The separate effects of target RGB, geometry and response inference remain
unmeasured under matched inputs.
\begin{table}[ht]
\centering\small\begin{tabular}{llr}
\toprule
Method & Origin & IoU \\
\midrule
NVOS \cite{ludvig} & P & 70.10 \\
SA3D-TRF \cite{ludvig} & P & 90.30 \\
SA3D-GS \cite{ludvig} & P & 92.20 \\
SAGA \cite{ludvig} & P & 92.60 \\
OmniSeg3D \cite{ludvig} & P & 91.70 \\
LUDVIG-DINOv2 \cite{ludvig} & P & 92.40 \\
LUDVIG-SAM \cite{ludvig} & P & 91.30 \\
LUDVIG-SAM2 \cite{ludvig} & P & 91.30 \\
VCAR \cite{vcar} & P & 93.50 \\
LEGO \cite{lego} & P & 94.20 \\
WildSeg3D \cite{wildseg3d} & P & 94.10 \\
Trace3D \cite{trace3d} & P & 92.50 \\
LUDVIG-SAM & R & 91.26 \\
Ours (v10) & O & 86.51 \\
Ours + coverage (N01) & M & 87.78 \\
Ours + coverage + source readout (N02) & M & 87.58 \\
Ours + local edit credit (N03) & M & 87.55 \\
Ours + image response + 3D (N04) & M & 88.32 \\
Ours + partial image response + 3D (N05) & M & 85.38 \\
Ours + unboxed image response + 3D (N06) & M & 84.15 \\
Ours + reference memory + 3D (N07) & M & 87.89 \\
Ours + native memory writeback (N08) & M & 88.11 \\
Ours + static native writeback (N08 secondary) & M & 88.84 \\
Ours + signed reference affinity (N09) & M & 81.58 \\
Ours + direct signed affinity (N10) & M & 77.73 \\
Ours + instance-constrained extent (N11) & M & 89.97 \\
Ours + signed reference completion (N12) & M & 91.48 \\
Ours + paired reference delta (N13) & M & 90.75 \\
Ours + source feasibility admission (N14) & M & 92.94 \\
Ours + source quality admission (N14 secondary) & M & 93.17 \\
Ours + completed multiview fusion (N15) & M & 93.03 \\
Ours + query-local target decoding (N16) & M & 93.28 \\
Ours + full-frame target decoding (N16 control) & M & 93.10 \\
Ours (v10, selected set) & O & 85.22 \\
\bottomrule
\end{tabular}

\caption{NVOS eight-task IoU. P values are reported in LUDVIG Tables 1/D~\cite{ludvig}.
R averages all three seeds in each task. B uses a different reference/source
contract; its new learned replacement remains pending.}
\label{tab:nvos}
\end{table}

\paragraph{Multimodal UQIS.}
UQIS is the local ScanNet-UQIS9 v0.3.1 protocol; no external published row in
the cited papers evaluates these same 268 queries. Table~\ref{tab:uqis} therefore
reports all modalities and a retained regional-readout baseline. The difference
between regional readout and v10 belongs to the preserved complete decoder,
not the new selector. The matched v10-to-replacement comparison isolates the
response decision. The full
shared policy improves image IoU by 1.088 points and AP by 1.268 points. One
query improves and none are harmed. Text and both actual-point outputs remain
unchanged; the point rows verify preservation rather than learned improvements.
All reported outputs use the same policy, without choosing a different winning
variant for text, images or categories.
\begin{table}[ht]
\centering\small\resizebox{\linewidth}{!}{\begin{tabular}{llrrrrrrrr}
\toprule
Method & Origin & T-IoU & T-AP & I-IoU & I-AP & P2-IoU & P2-AP & P3-IoU & P3-AP \\
\midrule
Ours (regional readout) & O & 12.11 & 23.69 & 26.55 & 37.75 & 42.09 & 52.74 & 41.52 & 52.36 \\
Ours (v10) & O & 16.39 & 26.60 & 39.90 & 59.31 & 49.74 & 75.79 & 49.05 & 75.33 \\
Ours + replacement (L11) & M & 16.39 & 26.60 & 40.99 & 60.58 & 49.74 & 75.79 & 49.05 & 75.33 \\
\bottomrule
\end{tabular}
}
\caption{UQIS matched comparison. T/I/P2/P3 denote text, image, 2D point and 3D
point. Actual selected-set IoU and continuous-response AP are distinct metrics.}
\label{tab:uqis}
\end{table}

\subsection{Mechanism ablations}
All rows in Table~\ref{tab:ablations} retain the same teachers, complete
candidates, supervision rooms and output definitions. The query-blind arm returns only the fitted structural prior (64 trees,
versus 128 in the full model); it therefore also changes capacity. Permutation
instead uses the fitted full model and
shuffles all teacher scores jointly within each view. It preserves their
multisets and correlations. Removing the null constraint retains the same
64+64-tree capacity and total fit weight, replacing permuted training examples
with duplicates of aligned examples. Each trained control uses the same
calibration procedure. The zero-margin control uses the full fitted model and
changes only the acceptance margin. These are response-level ablations.

\begin{table}[ht]
\centering\small\resizebox{\linewidth}{!}{\begin{tabular}{llrrrrrrr}
\toprule
Policy & Repair / harm & S19 & S15 & S10 & T-IoU & I-IoU & T-AP & I-AP \\
\midrule
Incumbent only & 0 / 0 & 61.01 & 62.53 & 72.49 & 16.39 & 39.90 & 26.60 & 59.31 \\
Full shared replacement & 1 / 0 & 61.01 & 62.53 & 72.49 & 16.39 & 40.99 & 26.60 & 60.58 \\
Without query coordinates & 1 / 1 & 61.01 & 62.53 & 72.49 & 16.39 & 40.26 & 26.60 & 59.50 \\
Permuted query correspondence & 0 / 0 & 61.01 & 62.53 & 72.49 & 16.39 & 39.90 & 26.60 & 59.31 \\
Without null correction constraint & 1 / 0 & 61.01 & 62.53 & 72.49 & 16.39 & 40.99 & 26.60 & 60.58 \\
Without positive margin & 56 / 33 & 60.54 & 62.05 & 71.62 & 22.17 & 41.17 & 35.18 & 60.53 \\
\bottomrule
\end{tabular}
}
\caption{Matched controls on all ten ScanNet rooms and 268 UQIS queries.
S19/15/10 are mesh mIoU. Repair/harm counts concern actual UQIS set IoU;
unchanged actual points are included in the denominator. This is one coupled
policy comparison, not a collection of per-benchmark best runs.}
\label{tab:ablations}
\end{table}

The image gain over query-blind and permuted controls is respectively 0.725 and
1.088 points; descriptive nine-room bootstrap intervals are [0,2.174] and
[0,3.264]. They do not establish a robust query-specific advantage. The full
and no-null arms make identical choices in all 144 text/image/menu decisions,
so the null constraint has no demonstrated accuracy contribution. Removing the
positive margin improves text and image IoU, but reduces ScanNet mIoU, including
a 0.869-point loss on the ten-class menu. Retention prevents some damage while
also rejecting useful instance replacements. Neither the full policy nor the
zero-margin variant is promoted to replace the reference system.

\subsection{Candidate capacity and failure attribution}
For a fixed candidate pool, let $G_0$ be gain under zero-margin selection and
$G_\tau$ gain under calibrated retention. The exact diagnostic is
\[
G_{\rm pool}=R_{\rm rank}+(G_0-G_\tau)+G_\tau.
\]
The middle term is signed: it is negative when retention prevents damage.
This identity prevents attributing all unused candidate capacity to the margin.
For text, 15.395 points of available IoU gain decompose into 9.615 ranking regret,
5.780 retention cost and zero actual gain. For images, 9.492 available points
decompose into 8.224 ranking regret, 0.180 retention cost and 1.088 actual gain.
Thus image ranking is the dominant measured bottleneck; lowering the margin
cannot recover most of the unused capacity. Category mean-menu capacity is
only 0.913 points, and its negative retention cost indicates prevented damage.
These label-oracle diagnostics are computed after prediction and are never
automatic method results.

\paragraph{Five-task ablation coverage.}
The same causal questions apply to all five outputs. Current coverage and the
minimal remaining tests are listed in Table~\ref{tab:coverage}. Historical full
cohort fusion/solver controls remain in the supplement; they cannot substitute
for the learned mechanism's missing LERF/NVOS experiments. No new systematic
feature-field ablation is required.
\begin{table}[ht]
\centering\small
\begin{tabular}{lp{.25\linewidth}p{.41\linewidth}}
\toprule
Task & Completed mechanism evidence & Remaining matched test \\
\midrule
LERF-2D & Provider and historical fusion controls & Complete-response alternatives; full / blind / permuted / zero-margin. \\
LERF-3D & Independent provider selection scored & Alternatives for the actual selection; same query decision across views. \\
ScanNet OVS & Full / blind / permuted / no-null / zero-margin & Cross-provider transfer with input contract fixed. \\
NVOS & Provider and historical solver controls & Reference/track evidence that works with one region per view; same decision controls. \\
UQIS & All five controls; 268 queries scored & Independent teacher configurations and unseen rooms; actual-point interventions if claimed. \\
\bottomrule
\end{tabular}
\caption{Required mechanism evidence. Planned cells are not completed ablations.
Teacher-transfer tests must hold each provider fixed within its paired comparison.}
\label{tab:coverage}
\end{table}
